"CRUD - Create, Read, Update and Delete"

---------------------------------------------------
del lista[2] 
-> Aqui o "del" é método que vai deletar;
-> "lista" é o nome da lista;
-> "[2]" é o item da lista que vai ser deletado.
---------------------------------------------------


---------------------------------------------------
lista.append(50)
-> Aqui "lista" é o nome da lista;
-> ".append" é o método que está sendo utilizado para adicionar à lista
-> "(50)" é o valor que será adicionado ao FINAL da lista.
---------------------------------------------------


---------------------------------------------------
lista.pop()
-> "lista" é o onome da lista;
-> ".pop()" é o método usado para remover algum item da lista.

OBS.: quando não é colocado nada dentro do parenteses, ele remove
sempre o último elemento da lista.

OBS.: já, para remover um item em específico da lista, eu preciso
colocar o índice dele dentro do parenteses.

OBS.: se eu fizer "v = ultimo_valor = lista.pop()
a variável 'v' vai me mostrar o elemento que foi removido da lista."
---------------------------------------------------


---------------------------------------------------
lista.insert(1, 'cla 200')
-> "lista" é o nome da lista;
-> "insert" é o método que estou utilizando para adicionar à lista;
-> "(1,...)" é o índice onde quero adicionar meu novo item;
-> "(..., 'cla 200')" é o item que estou adicionando à lista.
---------------------------------------------------


---------------------------------------------------
lista.clear()
-> Serve para limpar a lista inteira.
---------------------------------------------------


---------------------------------------------------
Para concatenar duas lista eu posso fazer de dois modos:

lista_a = [1, 2, 3]
lista_b = [4, 5, 6]

lista_a = lista_a + lista_b

isso vai me dar o mesmo resultado que:

lista_a.extend(lista_b)
---------------------------------------------------


---------------------------------------------------
Para copiar uma lista em outra lista:

lista_a = ['Luiz', 'Maria', 1, True, 1.2]
lista_b = lista_a.copy()

- Se eu fizer:

lista_a = ['Luiz', 'Maria', 1, True, 1.2]
lista_b = lista_a

Tudo que eu alterar em "lista_a" vai mudar em "lista_b"
e tudo que eu alterar em "lista_b" vai mudar em "lista_a".

Por isso é melhor usar o método ".copy()"
---------------------------------------------------


---------------------------------------------------
print(*lista)
Me tras tudo alinhado, como se tudos os itens da lista
estivessem em uma só linha.

print(*lista, sep = '\n')
Quando uma lista estiver dentro de outra, vai retornar
lista por lista, como se fosse um "for" dentro
de outro "for".
---------------------------------------------------


---------------------------------------------------
Listas:
lista_a = []
- É mutável, ou seja, posso modificá-la, adicionando itens, removando (CRUD);

Tuple;
tupla_a = 'item_1', 'item_2', 'item_3'
- É mais prático. Diferente da lista, a tupla eu uso quando não for auterar nada.
Ela se torna mais rápida e leve de se usar. 